\exercisesheader{}

% 17

\eoce{\qt{Konsumsi alkohol di bawah umur, Bagian I\label{underage_drinking_intro}}
Data yang dikumpulkan oleh Substance Abuse and Mental Health
Services Administration (SAMHSA) menunjukkan bahwa 69.7\% kaum muda
berusia 18--20 tahun mengonsumsi minuman beralkohol dalam suatu
tahun.\footfullcite{webpage:alcohol}
\begin{parts}
\item Misalkan diambil sampel acak berisi sepuluh orang berusia 18--20 tahun.
Apakah distribusi binomial tepat digunakan untuk menghitung peluang bahwa
tepat enam orang mengonsumsi minuman beralkohol? Jelaskan.
\item Hitung peluang bahwa tepat 6 dari 10 orang berusia 18--20 tahun yang
dipilih secara acak mengonsumsi minuman beralkohol.
\item Berapakah peluang bahwa tepat empat dari sepuluh orang berusia 18--20
tahun \textit{tidak} mengonsumsi minuman beralkohol?
\item Berapakah peluang bahwa paling banyak 2 dari 5 orang berusia 18--20
tahun yang dipilih secara acak mengonsumsi minuman beralkohol?
\item Berapakah peluang bahwa setidaknya 1 dari 5 orang berusia 18--20 tahun
yang dipilih secara acak mengonsumsi minuman beralkohol?
\end{parts}
}{}

% 18

\eoce{\qt{Cacar air, Bagian I\label{chicken_pox_intro}} Boston Children's
Hospital memperkirakan bahwa 90\% orang Amerika pernah terkena cacar air
sebelum mencapai usia dewasa. \footfullcite{bostonchildrenshospital:chickenpox}
\begin{parts}
\item Misalkan kita mengambil sampel acak berisi 100 orang dewasa Amerika.
Apakah distribusi binomial tepat digunakan untuk menghitung peluang bahwa
tepat 97 dari 100 orang dewasa Amerika yang dipilih secara acak pernah terkena
cacar air pada masa kanak-kanak? Jelaskan.
\item Hitung peluang bahwa tepat 97 dari 100 orang dewasa Amerika yang dipilih
secara acak pernah terkena cacar air pada masa kanak-kanak.
\item Berapakah peluang bahwa tepat 3 orang dari sampel baru yang berisi 100
orang dewasa Amerika \textit{tidak} pernah terkena cacar air pada masa
kanak-kanak?
\item Berapakah peluang bahwa setidaknya 1 dari 10 orang dewasa Amerika yang
dipilih secara acak pernah terkena cacar air?
\item Berapakah peluang bahwa paling banyak 3 dari 10 orang dewasa Amerika
yang dipilih secara acak \textit{tidak} pernah terkena cacar air?
\end{parts}
}{}

% 19

\eoce{\qt{Konsumsi alkohol di bawah umur, Bagian II\label{underage_drinking_normal_approx}}
Pada Latihan~\ref{underage_drinking_intro}, kita mengetahui bahwa sekitar
70\% orang berusia 18--20 tahun mengonsumsi minuman beralkohol dalam suatu
tahun. Sekarang kita mempertimbangkan sampel acak berisi lima puluh orang
berusia 18--20 tahun.
\begin{parts}
\item Berapa banyak orang yang diperkirakan telah mengonsumsi minuman
beralkohol? Berapakah simpangan bakunya?
\item Apakah Anda akan terkejut jika terdapat 45 orang atau lebih yang telah
mengonsumsi minuman beralkohol?
\item Berapakah peluang bahwa 45 orang atau lebih dalam sampel ini telah
mengonsumsi minuman beralkohol? Bagaimana peluang ini berkaitan dengan jawaban
Anda pada bagian (b)?
\end{parts}
}{}

% 20

\eoce{\qt{Cacar air, Bagian II\label{chicken_pox_normal_approx}} Pada
Latihan~\ref{chicken_pox_intro}, kita mengetahui bahwa sekitar 90\% orang
dewasa Amerika pernah terkena cacar air sebelum dewasa. Sekarang kita
mempertimbangkan sampel acak berisi 120 orang dewasa Amerika.
\begin{parts}
\item Berapa banyak orang dalam sampel ini yang diperkirakan pernah terkena
cacar air pada masa kanak-kanak? Berapakah simpangan bakunya?
\item Apakah Anda akan terkejut jika terdapat 105 orang yang pernah terkena
cacar air pada masa kanak-kanak?
\item Berapakah peluang bahwa 105 orang atau kurang dalam sampel ini pernah
terkena cacar air pada masa kanak-kanak? Bagaimana peluang ini berkaitan dengan
jawaban Anda pada bagian (b)?
\end{parts}
}{}

% 21

\eoce{\qt{Permainan dreidel\label{dreidel}} Dreidel adalah gasing bersisi empat
yang masing-masing sisinya memuat salah satu huruf Ibrani \textit{nun},
\textit{gimel}, \textit{hei}, dan \textit{shin}. Pada sekali putaran dreidel,
setiap sisi memiliki peluang yang sama untuk muncul. Misalkan Anda memutar
dreidel tiga kali. Hitung peluang memperoleh

\noindent\begin{minipage}[c]{0.45\textwidth}
\begin{parts}
\item setidaknya satu \textit{nun}?
\item tepat 2 \textit{nun}?
\item tepat 1 \textit{hei}?
\item paling banyak 2 \textit{gimel}? \vspace{3mm}
\end{parts}
\end{minipage}%
\begin{minipage}[c]{0.25\textwidth}
\ \vspace{2mm}

\Figures[Gambar dua dreidel kayu.]{0.95}{eoce/dreidel}{dreidel.jpg}\vspace{2mm}
\end{minipage}%
\begin{minipage}[c]{0.28\textwidth}%
{\footnotesize Foto oleh Staccabees, dipangkas \\
  (\oiRedirect{textbook-flickr_staccabees_dreidels}{http://flic.kr/p/7gLZTf}) \\
  \oiRedirect{textbook-CC_BY_2}{lisensi CC~BY~2.0}} \\
\end{minipage}
}{}
\vfill

% 22

\eoce{\qt{Araknofobia\label{arachnophobia}}
Sebuah jajak pendapat Gallup menemukan bahwa 7\% remaja (berusia 13 hingga
17 tahun) mengalami araknofobia dan sangat takut terhadap laba-laba. Di
sebuah perkemahan musim panas, terdapat 10 remaja yang tidur di setiap tenda.
Asumsikan bahwa 10 remaja ini saling independen.%
\footfullcite{webpage:spiders}
\begin{parts}
\item Hitung peluang bahwa setidaknya satu di antara mereka mengalami
araknofobia.
\item Hitung peluang bahwa tepat 2 di antara mereka mengalami araknofobia.
\item Hitung peluang bahwa paling banyak 1 di antara mereka mengalami
araknofobia.
\item Jika pembina perkemahan ingin memastikan bahwa tidak lebih dari 1
remaja di setiap tenda takut terhadap laba-laba, apakah menempatkan remaja ke
dalam tenda secara acak tampak masuk akal?
\end{parts}
}{}
\vfill

% 23

\eoce{\qt{Warna mata, Bagian II\label{eye_color_binomial}}
Latihan~\ref{eye_color_geometric} memperkenalkan pasangan suami istri bermata
cokelat yang memiliki peluang 0.75 untuk mempunyai anak bermata cokelat,
peluang 0.125 untuk mempunyai anak bermata biru, dan peluang 0.125 untuk
mempunyai anak bermata hijau.
\begin{parts}
\item Berapakah peluang bahwa anak pertama mereka bermata hijau dan anak
kedua tidak bermata hijau?
\item Berapakah peluang bahwa tepat satu dari kedua anak mereka bermata hijau?
\item Jika mereka mempunyai enam anak, berapakah peluang bahwa tepat dua di
antaranya bermata hijau?
\item Jika mereka mempunyai enam anak, berapakah peluang bahwa setidaknya
satu di antaranya bermata hijau?
\item Berapakah peluang bahwa anak pertama yang bermata hijau adalah anak
ke-$4$?
\item Apakah hanya 2 dari 6 anak mereka yang bermata cokelat akan dianggap
tidak biasa?
\end{parts}
}{}
\vfill

% 24

\D{\newpage}

\eoce{\qt{Anemia sel sabit\label{sickle_cell_anemia}} Anemia sel sabit adalah
kelainan darah genetik yang membuat sel darah merah kehilangan kelenturannya
dan mengambil bentuk ``sabit" yang abnormal dan kaku, sehingga menimbulkan
risiko berbagai komplikasi. Jika kedua orang tua adalah pembawa penyakit ini,
seorang anak memiliki peluang 25\% untuk menderita penyakit tersebut, peluang
50\% untuk menjadi pembawa, dan peluang 25\% untuk tidak menderita penyakit
tersebut maupun menjadi pembawa. Jika dua orang tua yang merupakan pembawa
penyakit ini mempunyai 3 anak, berapakah peluang bahwa
\begin{parts}
\item dua anak akan menderita penyakit tersebut?
\item tidak satu pun anak akan menderita penyakit tersebut?
\item setidaknya satu anak tidak akan menderita penyakit tersebut maupun
menjadi pembawa?
\item anak pertama yang menderita penyakit tersebut merupakan anak ke-$3$?
\end{parts}
}{}
\vfill

% 25

\eoce{\qt{Menjelajahi permutasi\label{explore_combinations}} Rumus untuk
banyaknya cara menyusun $n$ objek adalah $n! = n\times(n-1)\times \cdots
\times 2 \times 1$. Latihan ini memandu Anda menurunkan rumus tersebut untuk
beberapa kasus khusus.

\indent Sebuah perusahaan kecil memiliki lima karyawan: Anna, Ben, Carl,
Damian, dan Eddy. Perusahaan tersebut memiliki lima tempat parkir yang
berderet, tidak ada yang ditetapkan bagi karyawan tertentu, dan setiap hari
para karyawan memarkir mobil di tempat yang dipilih secara acak. Dengan kata
lain, semua kemungkinan urutan mobil dalam deretan tempat parkir tersebut
memiliki peluang yang sama.
\begin{parts}
\item Pada suatu hari, berapakah peluang bahwa para karyawan memarkir mobil
dalam urutan abjad?
\item Jika urutan abjad memiliki peluang yang sama untuk terjadi seperti
setiap urutan lain yang mungkin, ada berapa cara menyusun kelima mobil itu?
\item Sekarang, pertimbangkan sampel berisi 8 karyawan. Ada berapa cara yang
mungkin untuk mengurutkan mobil milik 8 karyawan tersebut?
\end{parts}
}{}
\vfill

% 26

\eoce{\qt{Anak laki-laki\label{male_children}} Walaupun peluang memiliki anak
laki-laki atau anak perempuan sering dianggap sama, peluang sebenarnya untuk
memiliki anak laki-laki sedikit lebih tinggi, yaitu 0.51. Misalkan sepasang
suami istri berencana mempunyai 3 anak.
\begin{parts}
\item Gunakan model binomial untuk menghitung peluang bahwa dua di antaranya
adalah anak laki-laki.
\item Daftarkan semua kemungkinan urutan 3 anak jika 2 di antaranya adalah
anak laki-laki. Gunakan aturan penjumlahan untuk hasil-hasil yang saling lepas
guna menghitung peluang pada bagian (a), lalu pastikan kedua jawaban sama.
\item Untuk peluang bahwa, dari 8 anak mereka, 3 adalah anak laki-laki,
jelaskan secara singkat mengapa pendekatan pada bagian (b) lebih merepotkan
daripada pendekatan pada bagian (a).
\end{parts}
}{}
\vfill
